\section{Installation}

\subsection{Scope of Delivery}
\begin{itemize}
	\item 	LARUS Vario Display
	\item 	Mounting screws
	\item 	1:1 standard RJ45 cable
	\item 	Micro SD card with adapter
	\item 	D-Sub-9 solder connector and housing
	\item 	1,5 mm HEX key for mounting the rotary knobs
\end{itemize}

\subsection{Design and Function}

LARUS Vario Display shows the data measured and calculated by LARUS. LARUS is an advanced
variometer with real time wind measurement capabilities. It incorporates state of the art pressure
sensors, an advanced IMU and GNSS receivers to gather precise flight information data.

Key features of the display are:
\begin{itemize}
	\item 	Round display for 57 mm standard instrument panel cutouts
	\item 	Bright and colourful screen
	\item 	Lightweight, compact design with black anodised aluminium housing
	\item 	Two rotary knobs with push-button function to change settings and enter the menus
1,5 mm HEX key for mounting the rotary knobs
\end{itemize}

LARUS Vario Display is designed and continuously improved by Prof. Dr. Klaus Schaefer, Maximilian
Betz, Winfried Simon, Peter Simon and the SteFly team.


\subsection{System Architecture}


\subsubsection{Single Seater Configuration}
\includegraphics[width=16cm]{config-singleseater.png}

\subsubsection{Double Seater Configuration}
\includegraphics[width=16cm]{config-doubleseater.png}


\subsection{Connectors and Cabling}

\subsubsection{Overview of the Connectors}
\includegraphics[width=8cm]{connectors-overview.png}

\subsubsection{CAN and RS232 Ports}
\includegraphics[width=8cm]{connectors-rj45.pdf}

\begin{table*}[ht!]
	\begin{tabular}{C{0.1\linewidth} L{0.40\linewidth} L{0.40\linewidth}}
		\toprule
		\textbf{Pin} & \textbf{CAN} & \textbf{RS232} \\
		\midrule
		1 & GND (internally connected) & GND (internally connected) \\
		2 & GND (internally connected) & GND (internally connected) \\
		3 & NC 						   & RS232-1-RX \\
		4 & CAN Low                    & RS232-1-TX \\
		5 & CAN High                   & NC \\
		6 & NC                         & NC \\
		7 & VCC [9-28V DC]             & VCC [9-28V DC] \\
		8 & VCC [9-28V DC]             & VCC [9-28V DC] \\
		\bottomrule
	\end{tabular}
	\caption{Pin Asssignment CAN and RS232 Rj45}
\end{table*}

\subsubsection{GPIO / D-SUB 9}
Several other switches, sensors and devices can be connected via the d-sub connector. Please note
that not all functions have been implemented at date of release of this manual.

The following sketch shows the view into the male connector of the LARUS Vario Display.
\includegraphics[width=6cm]{connector-dsub9.png}

\begin{table*}[ht!]
	\begin{tabular}{C{0.1\linewidth} L{0.40\linewidth} L{0.20\linewidth}}
		\toprule
		\textbf{Pin} & \textbf{CAN} & \textbf{RS232} \\
		\midrule
		1 & GND  					& (Ground) \\
		2 & DI3 - Gear 				& Input \\
		3 & DI1 - Water Ballast		& Input \\
		4 & DO2						& Output \\
		5 & GND						& (Ground) \\
		6 & DI4 - Speed Breakes 	& Input \\
		7 & DI2 - Spped to fly		& Input \\
		8 & GND						& (Ground) \\
		9 & DO1 - Canopy Flasher 	& Output \\
		\bottomrule
	\end{tabular}
	\caption{Pin Asssignment GPIO / D-SUB 9}
\end{table*}

For easier pin identification these numbers are also moulded into the female connector (scope of
delivery).

After wiring, settings must be made in the LARUS Vario Display.


\pagebreak

\subsubsection{Audio}
An audio socket is available for connecting a loudspeaker with a 3.5 mm jack plug. The internal
resistance of the speaker must be between 4-8 \si{\ohm} (max. output of 3 W @ 8 \si{\ohm}).

\picnote
	{pictograph-yellow-warning.pdf}
	{A single loudspeaker must not be connected to more than one device.}

\subsubsection{SD-Card}
The device has an SD card slot for firmware updates.

\picnote
	{pictograph-yellow-warning.pdf}
	{As SD card extensions can damage the LARUS Vario Display, we do not accept any
	liability for damage resulting from their use.}

\subsubsection{Cabling}

If a serial connection (RS232) is established between SteFly NAV and LARUS Vario Display, the MC
value settings, for example, is synchronised/transferred between the devices. In addition, the LARUS
Vario Display processes inputs from a Speed to Fly / Vario switch that is connected directly to the
SteFly Remote Control Stick.
Thefollowing steps are required:
\begin{itemize}
	\item 	Connect devices with cable:
	\begin{itemize}
		\item 	The LARUS Sensor Box CAN-port needs to be connected with 1:1 patch cables with
				the CAN-port of the LARUS VARIO Display (in double seater configurations using the
				cablebox)
		\item 	The devices running XCSoar can be connected between Vario Display RS232 port and
				SteFly NAV ttySx port with patch cable RX/TX crossed
	\end{itemize}
	\item 	Optionally connect switches for automatic adjustment of water ballast and gear warning to
			the appropriate DSUB-9 pin and any GND pin. For the gear warning you have the option to
			use two distinct directly connected switches for the gear and speed brake (recommended) or
			one cable with the two switches connected in series to the DSUB-9 connector.
	\item 	The SteFly Remote Control Stick can be connected in two different ways:
	\begin{itemize}
		\item 	if available: with a separate cable to the DSUB-9 pin 7
		\item 	Using the USB-Cable (always available)
	\end{itemize}
	\item	Adjust settings in XCSoar (Version 7.44 or higher) / OpenSoar (7.43 or higher):
	\begin{itemize}
		\item 	If the Vario Display is connected directly to the SteFly NAV, select the corresponding
				ttyS port: Config -> Devices -> select a free line and enter the port, Baud rate 38400,
				Driver Larus, Sync. to device on
	\end{itemize}
	\item 	Adjust settings in LARUS Vario Display settings for the connected optional switches
	\begin{itemize}
		\item 	Gear warning
		\item 	water ballast
		\item 	StF Switch configuration (if not automatically switched)
		\begin{itemize}
			\item 	Input Pin: To use an external switch (at the control stick, flaps lever etc), set
					Vario Control to Input Pin and choose the right setting in StF Pin Config
					(active: when open/closed)
			\item 	NMEA for SteFly Remote Control Stick. In addition to the setting in the vario
					you need to enable install an event definition file in XCSoar/OpenSoar
					(Config / System / Look / Language, Input / Events – check Expert, click on
					Events – Download – GLB-XCI-xcremote-XCNAV.xci). Leave XCSoar config and
					restart XCSoar.
		\end{itemize}
	\end{itemize}
\end{itemize}


\subsection{CAN Termination}

LARUS Vario Display and LARUS Sensor Box are connected via CAN. CAN bus networks require
termination resistors at each end of the network. 

Therefore, all units have integrated a switch to activate the resistor:

\includegraphics[width=16cm]{can-termination.png}

\begin{table*}[ht!]
	\begin{tabular}{L{0.25\linewidth} C{0.156\linewidth} C{0.156\linewidth} C{0.156\linewidth} C{0.156\linewidth}}
		\toprule
		\textbf{Description} & \textbf{Display front} & \textbf{Larus Box} & \textbf{CAN Splitter} & \textbf{Display rear} \\
		\midrule
		Single Seater 							& on  & on  & -   & -   \\
		Double Seater, Larus box in front panel & off & on  & off & on  \\
		Double Seater, Larus box in rear panel  & on  & on  & off & off \\
		\bottomrule
	\end{tabular}
	\caption{CAN Termination Single / Double Seater}
\end{table*}

Please note: All LARUS Sensor Boxes delivered before March 2025 have no CAN termination switch.
CAN termination is always on from default.



\subsection{External Fuse}

\picnote
	{pictograph-yellow-warning.pdf}
	{Usually, LARUS Vario Display is powered by LARUS via patch cable between the CAN
	ports. LARUS needs to be protected by an external fuse (500 mA to max. 3A) like it is
	common practice for all electric devices in aviation. If LARUS gets its energy from
	another main instrument (e.g. SteFly NAV via D-Sub connector) please make sure
	that the main instrument is protected by an external fuse.}

\subsection{Installation Location}

The following picture shows a typical installation situation of LARUS Vario Display in the instrument
panel of a glider.

\includegraphics[width=8cm]{panel.png}

The display fits in a standard 57 mm cutout and is fastened by 3 M3 screws.

\includegraphics[width=8cm]{knobs-and-screws.png}

For installation it is required to remove the two rotary knobs with a HEX 1,5mm key.


\subsection{Installation Orientation}

\picnote
	{pictograph-yellow-warning.pdf}
	{If you intend to replace an existing instrument with the Lars Vario Display, check the
	desired installation position before drilling the 7.3 mm hole for the rotary encoder due
	to the asymmetrical housing to the display.}

The display can be mounted in the orientations 0° / 90° / 180° / 270°.

After mounting the display, its orientation needs to be configured in the menu Display Rotation.

	

\subsection{Initial Operation and Function Test}

For initial operation please follow these steps:
\begin{enumerate}
	\item 	Check that the LARUS Vario Display is connected to LARUS Sensor Box or CAN Splitter
			according to the drawings in chapter 4.2
	\item 	Power on LARUS
	\item	Check that LARUS Vario Displays boots and a yellow or green satellite pictogram appears. The
			vario pointers should perform slight movements around zero position.
\end{enumerate}

\subsection{Maintenance}
The whole system has no serviceable parts. 

To obtain warranty service, please contact SteFly directly.

\picnote
	{pictograph-yellow-warning.pdf}
	{Opening the housing of LARUS Vario Display will void the warranty!}

\subsection{Troubleshooting}
\begin{table*}[ht!]
	\begin{tabular}{L{0.29\linewidth} L{0.29\linewidth} L{0.29\linewidth}}
		\toprule
		\textbf{Problem} & \textbf{Possible Causes} & \textbf{Solution} \\
		\midrule
		LARUS Vario Display boots, but satellite pictogram is red and vario pointers are fixed
		& Connection to LARUS CAN port with Rx / Tx crossed patch cable instead of standard 1:1 patch cable
		& Change the patch cable and use the one supplied in scope of delivery \\
		LARUS Vario Display boots, but satellite pictogram is red and vario pointers are fixed
		& Connection to LARUS via RS232 port
		& Connect the CAN ports instead \\
		\bottomrule
	\end{tabular}
	\caption{Troubleshooting}
\end{table*}

\subsection{Firmware Update}

The LARUS team is continuously improving the software and releasing updates e.g. on LARUS Vario
Display web page.

To update the firmware, put the new .bin-file on the SD card included in scope of delivery and insert
it into the SD-card slot of LARUS Vario Display. The device will only install compatible firmware
releases greater than the release installed. After powering on, the device is booting within less than a
second. If a compatible firmware is detected on the SD Card, the display keeps black for about 3-5
seconds before displaying the a Firmware Upgrade is in process. Afterwards the name of the new
firmware version will appear.

If the update is not initialised after powering on the device with an inserted SD card containing a new
firmware release, please use a different SD card.

Requirements for the SD card: min. 4 GB SDHC, FAT32 format (not GPT)


